%%%%%%%%%%%%Outline%%%%%%%%%%%%%%%%%%

%%%
% message<xx>
% xx

%%%%%%%%%%%%Outline%%%%%%%%%%%%%%%%%%

\section{Background and Motivation}
\label{sec:background-and-motivation}

\subsection{eBPF Architecture}

eBPF extends the original Berkeley Packet Filter~\cite{ebpf_sosp}
with a register-based instruction set, a rich type system, and a
general-purpose execution environment inside the Linux
kernel~\cite{bpf_lwn}. An eBPF program is compiled from C (or Rust)
into BPF bytecode by Clang/LLVM, producing an ELF object file. At
load time, the kernel subjects the bytecode to a static verifier
(\texttt{bpf\_check()}) that enforces memory safety, bounded
execution, and type correctness. Programs that pass verification are
JIT-compiled into native machine code and attached to kernel hooks
(XDP, TC, tracepoints, kprobes, LSM, \texttt{sched\_ext}, etc.).

The eBPF runtime provides three mechanisms for programs to interact
with kernel state: \emph{maps} (shared key-value data structures),
\emph{helper functions} (kernel-provided utility functions callable
from BPF), and \emph{kfuncs} (typed kernel functions exposed via BTF
that can be called like regular functions). Kfuncs are a recent
extension that allows kernel modules to export functions to BPF
programs with full type information, verified by the existing
\texttt{check\_kfunc\_call()} path without additional verifier
changes.

\subsection{Current JIT Limitations}

The kernel JIT compiler maps each BPF instruction to a single,
hand-audited native instruction sequence~\cite{linux_bpf_jit}. This
one-to-one mapping ensures that the generated code is safe and
auditable, but it also means that every instruction-selection decision
is fixed at kernel compile time. Three categories of optimization
opportunities are systematically missed:

\textbf{Absent platform-specific instructions.} The BPF instruction
set has no opcodes for conditional moves (\texttt{cmov}), bit-field
extraction (\texttt{bextr}), rotate (\texttt{rorx}), or
load-effective-address (\texttt{lea}). The kernel JIT therefore cannot
emit these instructions even when the BPF bytecode contains idioms
(e.g., shift-or pairs for rotation, branch-over-move diamonds for
conditional selection) that map directly to them.

\textbf{No runtime-guided optimization.} The kernel JIT performs a
single compilation pass with no feedback from runtime behavior. There
is no profile-guided optimization, no dynamic recompilation based on
hot-path information, no constant propagation from frozen map values,
and no workload-adaptive branch layout. Mature runtimes like JVM
HotSpot~\cite{jvm_hotspot} routinely employ all of these techniques.

\textbf{No extensible compilation framework.} Adding a new
optimization to the kernel JIT requires modifying core kernel source
files, passing multi-month upstream review, and waiting for the next
kernel release cycle. There is no plugin or pass infrastructure that
would allow new optimizations to be developed, tested, and deployed
independently of the kernel release schedule.

\subsection{Why Existing Approaches Fall Short}

Several systems have addressed BPF optimization, but none provides a
post-load, transparent, runtime-guided compilation framework
(Table~\ref{tab:existing_work}).

K2~\cite{k2} applies stochastic superoptimization to BPF bytecode.
Merlin~\cite{merlin} lifts BPF to LLVM IR for optimization.
EPSO~\cite{epso} uses evolutionary program synthesis at the bytecode
level. All three operate \emph{before} JIT compilation: they optimize
the bytecode input to the JIT, not its native output. Since the BPF
ISA lacks opcodes for \texttt{cmov}, \texttt{rorx}, or wide-load
fusion, no bytecode rewrite can introduce these native instructions.
Moreover, all three require the original \texttt{.bpf.o} file, cannot
optimize programs already loaded by third-party tools, and do not
support runtime profiling feedback.

ePass~\cite{epass} proposes in-kernel optimization passes at load
time, but requires source modifications and does not support post-load
re-optimization. BCF~\cite{bcf} offloads verification to userspace
via certificates, addressing verifier scalability rather than JIT
optimization. Neither system provides a runtime-guided, transparent
optimization framework.

\begin{table}[t]
  \centering
  \caption{Comparison with BPF optimization and verification systems.
  \textsc{BpfReJIT} is the only system that operates post-load,
  transparently, and with runtime profiling support.}
  \label{tab:existing_work}
  \small
  \begin{tabular}{lccccc}
    \toprule
    System & Timing & Transparent? & PGO? & Backend? & Safety \\
    \midrule
    K2~\cite{k2}          & pre-load  & No  & No  & No  & SMT solver \\
    Merlin~\cite{merlin}  & pre-load  & No  & No  & No  & compiler \\
    EPSO~\cite{epso}      & pre-load  & No  & No  & No  & compiler \\
    ePass~\cite{epass}    & load-time & No  & No  & No  & in-kernel \\
    BCF~\cite{bcf}        & load-time & No  & No  & No  & certificate \\
    \textbf{BpfReJIT}     & \textbf{post-load} & \textbf{Yes} & \textbf{Yes} & \textbf{Yes} & \textbf{verifier} \\
    \bottomrule
  \end{tabular}
\end{table}

\subsection{Quantifying the Code Quality Gap}
\label{sec:characterization}

\textbf{Methodology.} We compare the kernel JIT (Linux 6.15.11) against
llvmbpf on 31 pure-JIT microbenchmarks that isolate code generation
quality (no map lookups on the hot path). Each benchmark runs 30
iterations of 1{,}000 repeats; we report the per-benchmark median of
per-iteration medians. We use paired Wilcoxon signed-rank tests with
Benjamini-Hochberg correction and bootstrap 95\% CIs for suite-level
geomeans. Platform: Intel Core Ultra 9 285K (Arrow Lake-S), performance
governor, turbo disabled.

\begin{table}[t]
  \centering
  \caption{Kernel JIT code quality gap: suite-level summary.}
  \label{tab:exec_summary}
  \begin{tabular}{lrrr}
    \toprule
    Suite & Geomean (L/K) & 95\% CI & Sig.\ / Total \\
    \midrule
    Pure-JIT (31) & 0.849$\times$ & [0.834, 0.865] & 25/31 \\
    Runtime (9)   & 0.829$\times$ & [0.801, 0.859] & 7/9   \\
    \bottomrule
  \end{tabular}
\end{table}

Table~\ref{tab:exec_summary} summarizes the gap. llvmbpf is faster on
21/31 benchmarks; the kernel wins on 10/31, primarily loop-dominated
workloads where simpler code has lower I-cache pressure. Code size tells
a cleaner story: llvmbpf is smaller on \emph{all} 31 benchmarks (geomean
0.496$\times$), confirming the gap is systematic.

\textbf{Root causes.} JIT-dump analysis of 22 benchmarks decomposes the
kernel's instruction surplus into three mechanisms covering 89\% of the
total: byte-load recomposition (50.7\%), absent conditional moves
(19.9\%), and fixed callee-saved registers (18.5\%).

\subsection{The Case Against Fixed Kernel Heuristics}
\label{sec:policy_sensitivity}

A fixed kernel heuristic would close the gap only if every optimization
were unconditionally profitable. We present two forms of evidence that
this assumption fails.

\textbf{Evidence 1: Policy sensitivity.}
Table~\ref{tab:cmov_ablation} shows an ablation study in llvmbpf where
\texttt{cmov} emission is toggled. Disabling \texttt{cmov}
\emph{hurts} programs with unpredictable branches
(\texttt{switch\_dispatch}: +26.3\%) but \emph{helps} programs with
predictable branches (\texttt{bounds\_ladder}: $-$18.1\%). The same
legal lowering helps one workload class and hurts another by comparable
magnitudes. We validate both directions in the kernel JIT via BpfReJIT: forcing
CMOV on predictable-branch \texttt{log2\_fold} degrades performance by
34.5\% (383.5$\to$590.0\,ns, $p{=}0.002$), while applying it to
unpredictable-branch \texttt{cmov\_select} yields a modest improvement
(744$\to$723\,ns, \S\ref{sec:evaluation}).

\begin{table}[t]
  \centering
  \caption{Cmov ablation in llvmbpf: profitability depends on workload.
  Validated in the kernel JIT via BpfReJIT re-JIT (\S\ref{sec:evaluation}).}
  \label{tab:cmov_ablation}
  \begin{tabular}{lrrr}
    \toprule
    Benchmark & Normal & No-cmov & $\Delta$ \\
    \midrule
    \texttt{switch\_dispatch}  & 213\,ns & 269\,ns & \textbf{+26.3\%} \\
    \texttt{binary\_search}    & 205\,ns & 229\,ns & \textbf{+11.7\%} \\
    \texttt{bounds\_ladder}    & 83\,ns  & 68\,ns  & $-18.1\%$ \\
    \texttt{large\_mixed\_500} & 415\,ns & 316\,ns & $-23.9\%$ \\
    \bottomrule
  \end{tabular}
\end{table}

\textbf{Why static analysis is insufficient.} Branch predictability
depends on \emph{runtime input distributions}, not just program
structure. A BPF program performing packet classification against a
hash table has data-dependent branches whose predictability varies with
traffic mix---a property that changes between deployments and over time.
No static analysis of BPF bytecode can determine the distribution of
live traffic at a particular deployment site.

\textbf{Evidence 2: Combination interference.}
Even if each optimization were individually safe, applying all of them
simultaneously can produce regressions. We built a ``fixed-all'' kernel
with all four optimizations hardcoded. While code-size effects compose
additively (each program's native code shrinks by exactly the sum of the
individual deltas), execution-time effects \emph{do not}.
Table~\ref{tab:fixed_preview} previews the key finding: the fixed-all
kernel regresses \texttt{log2\_fold} by 28.3\% and
\texttt{stride\_load\_16} by 14.2\%.

\begin{table}[t]
  \centering
  \caption{Fixed-all kernel baseline: code size composes additively,
  execution time does not.}
  \label{tab:fixed_preview}
  \small
  \begin{tabular}{lrrr}
    \toprule
    Program & Exec $\Delta$ & Code $\Delta$ & Interpretation \\
    \midrule
    \texttt{rotate64\_hash}     & $-$16.8\% & $-$1{,}178\,B & Substrate win \\
    \texttt{pkt\_rss\_hash}     & $-$14.9\% & $-$151\,B  & Substrate win \\
    \texttt{load\_byte\_recomp} & $-$13.0\% & $-$12\,B   & Substrate win \\
    \texttt{stride\_load\_16}   & \textbf{+14.2\%} & $-$30\,B   & Combination regression \\
    \texttt{log2\_fold}         & \textbf{+28.3\%} & $-$2\,B    & CMOV regression \\
    \bottomrule
  \end{tabular}
\end{table}

The mechanism we need is not a better heuristic but a separation of
concerns: the kernel should enforce which native code variants are
\emph{safe to emit}, while the operator---who understands the
workload---specifies which safe variants to \emph{apply}, selectively
per program and per site.